\documentclass[12pt, a4paper]{report}

% --- Pacchetti Fondamentali ---
\usepackage[utf8]{inputenc}
\usepackage{lmodern}
\usepackage{geometry}
\geometry{a4paper, top=30mm, bottom=30mm, left=35mm, right=30mm}
\usepackage{graphicx}
\usepackage{amsmath}
\usepackage{amssymb}
\usepackage{booktabs}
\usepackage{xcolor}
\usepackage{listings}
\usepackage{hyperref}
\usepackage{fancyhdr}

% --- Configurazione Tracciamento Codice (Listings) ---
\definecolor{codegreen}{rgb}{0,0.6,0}
\definecolor{codegray}{rgb}{0.5,0.5,0.5}
\definecolor{codepurple}{rgb}{0.58,0,0.82}
\definecolor{backcolour}{rgb}{0.95,0.95,0.92}

\lstdefinestyle{codeStyle}{
    backgroundcolor=\color{backcolour},   
    commentstyle=\color{codegreen},
    keywordstyle=\color{magenta},
    numberstyle=\tiny\color{codegray},
    stringstyle=\color{codepurple},
    basicstyle=\ttfamily\footnotesize,
    breakatwhitespace=false,         
    breaklines=true,                 
    captionpos=b,                    
    keepspaces=true,                 
    numbers=left,                    
    numbersep=5pt,                  
    showspaces=false,                
    showstringspaces=false,
    showtabs=false,                  
    tabsize=2
}
\lstset{style=codeStyle}

% --- Configurazione Layout ---
\pagestyle{fancy}
\fancyhf{}
\rhead{\small\leftmark}
\lfoot{\small Bozza Tesi Triennale --- Brando Calderara}
\rfoot{\thepage}
\renewcommand{\headrulewidth}{0.4pt}
\renewcommand{\footrulewidth}{0.4pt}

\begin{document}

% ====================================================================
% FRONTESPIZIO
% ====================================================================
\begin{titlepage}
    \begin{center}
        \LARGE{\textbf{UNIVERSIT\`A DEGLI STUDI DI BRESCIA}}\\
        \vspace{5mm}
        \Large{Dipartimento di Ingegneria dell'Informazione}\\
        \vspace{3mm}
        \normalsize{Corso di Laurea in Ingegneria delle Tecnologie per l'Impresa Digitale (ITID)}\\
        
        \vspace{25mm}
        \textsl{\Large{Tesi di Laurea Triennale}}\\
        \vspace{10mm}
        \Huge{\textbf{Automazione e Centralizzazione del Monitoraggio di Infrastrutture Smart City via Zabbix API e Integrazioni Crittografiche}}\\
        \vspace{5mm}
        \Large{\textit{Caso di studio operativo presso Secoval S.r.l.}}
        
        \vspace{35mm}
        \begin{table}[h]
            \centering
            \normalsize
            \begin{tabular}{l p{3cm} l}
                \textbf{Relatore:} & & \textbf{Candidato:} \\
                Prof. Francesco Gringoli & & Brando Calderara \\
                & & Matr. [Inserire Matricola] \\
                \vspace{5mm}\\
                \textbf{Correlatore Aziendale:} & & \\
                [Nome Correlatore Secoval] & & 
            \end{tabular}
        \end{table}
        
        \vfill
        \begin{center}
            \line(1,0){350}\\
            \vspace{3mm}
            \normalsize{Anno Accademico 2025/2026}
        \end{center}
    \end{center}
\end{titlepage}

\clearpage

% ====================================================================
% INDICE
% ====================================================================
\tableofcontents
\clearpage

% ====================================================================
% INTRODUZIONE
% ====================================================================
\chapter*{Introduzione}
\addcontentsline{toc}{chapter}{Introduzione}
Il rapido sviluppo del paradigma delle \textbf{Smart City} ha portato le pubbliche amministrazioni locali a dotarsi di infrastrutture tecnologiche capillari e complesse. Reti di videosorveglianza urbana (TVCC), access point per il Wi-Fi pubblico, switch industriali di distribuzione e gateway LTE costituiscono oggi la spina dorsale dei servizi al cittadino. Gestire l'elevata eterogeneit\`a e la dislocazione geografica di tali dispositivi rappresenta una sfida cruciale per i provider di servizi tecnologici territoriali.

L'obiettivo di questo lavoro di tesi, svolto in collaborazione con l'azienda \textbf{Secoval S.r.l.}, \`e la progettazione e l'implementazione di un sistema di monitoraggio centralizzato, proattivo e automatizzato basato sulla piattaforma open-source \textbf{Zabbix 7.0}. 

Tradizionalmente, il monitoraggio di tali reti si \`e basato su controlli statici o su piattaforme commerciali limitate (come PRTG), che richiedevano un pesante overhead di gestione manuale per l'inserimento e la manutenzione dei dispositivi. Il nucleo innovativo di questo progetto risiede nell'estrazione automatica delle informazioni da sorgenti documentali aziendali (\textit{Configuration Management Database}), nel provisioning dinamico degli apparati tramite Zabbix API, nell'intercettazione in tempo reale di eventi di sicurezza sulle telecamere tramite protocolli nativi e nell'integrazione sicura (tramite crittografia asimmetrica e token JWT) con la piattaforma cloud di ticketing aziendale per la gestione dei guasti sul campo.

L'elaborato \`e strutturato come segue: il Capitolo 1 inquadra l'infrastruttura di rete di Secoval e le motivazioni tecnologiche della scelta di Zabbix 7.0. Il Capitolo 2 approfondisce le logiche di automazione dell'inventario tramite script Python e Zabbix API. Il Capitolo 3 tratta la configurazione dei controlli attivi e la gestione degli eventi TVCC. Il Capitolo 4 descrive l'architettura crittografica di interfacciamento con il sistema di ticketing, mentre il Capitolo 5 espone i meccanismi di reportistica KPI e manutenzione del sistema operati in ambiente Linux.

\clearpage

% ====================================================================
% CAPITOLO 1
% ====================================================================
\chapter{Il Contesto Operativo e Tecnologico}
\section{La gestione infrastrutturale in Secoval S.r.l.}
Secoval S.r.l. si occupa della gestione dei servizi tecnologici e informatici per conto di numerosi enti locali e comunit\`a montane. La sua infrastruttura di rete si estende su molteplici comuni, integrando dispositivi dedicati alla sicurezza stradale, alla connettivit\`a pubblica e al controllo del territorio.

\section{Eterogeneit\`a degli apparati monitorati}
La complessit\`a della rete Smart City monitorata deriva dalla convivenza di diverse tipologie di nodi critici:
\begin{itemize}
    \item \textbf{Telecamere di Videosorveglianza Urbana (TVCC):} Dispositivi IP (prevalentemente Hikvision) che richiedono monitoraggio sia di raggiungibilit\`a di rete che di integrit\`a del flusso video.
    \item \textbf{Access Point (AP):} Dispositivi per l'erogazione di connettivit\`a Wi-Fi pubblica nei centri urbani.
    \item \textbf{Switch di Rete:} Apparati industriali (L2/L3) che aggregano i flussi provenienti dalle periferiche urbane.
    \item \textbf{Dispositivi LTE/5G:} Router e gateway utilizzati per l'uplink in zone non raggiunte dalla fibra ottica, caratterizzati da indirizzamenti IP dinamici o connettivit\`a soggetta a fluttuazioni.
\end{itemize}

\section{Evoluzione dei Sistemi di Network Monitoring: da PRTG a Zabbix 7.0}
I limiti delle soluzioni precedentemente adottate includevano vincoli di licenziamento sul numero di sensori, scarsa attitudine all'automazione tramite script esterni e difficolt\`a nel gestire variazioni dinamiche degli host. La migrazione a \textbf{Zabbix 7.0 LTS} \`e stata dettata dalla sua architettura modulare, dal supporto nativo alle espressioni regolari per il log parsing, dalla presenza di una robusta interfaccia JSON-RPC (API) e dall'efficienza del demone \textit{Zabbix Proxy} per la raccolta dati distribuita sul territorio.

\clearpage

% ====================================================================
% CAPITOLO 2
% ====================================================================
\chapter{Automazione dell'Inventario e Configuration Management}
\section{Il foglio Excel aziendale come Source of Truth}
Per superare i limiti dei tradizionali meccanismi di \textit{Network Discovery} basati puramente su scansioni ICMP periodiche --- che rilevano gli host ma non sono in grado di assegnarvi metadati aziendali significativi --- \`e stato stabilito che il foglio Excel di gestione (CMDB aziendale) fungesse da unica sorgente di verit\`a. Ogni foglio di calcolo corrisponde a un comune e mappa in modo preciso campi quali: \texttt{ID Host}, \texttt{Indirizzo IP}, \texttt{ID Palo} (identificativo del posizionamento fisico) e le coordinate geografiche.

\section{Sviluppo del modulo di provisioning via Zabbix API}
\`E stato sviluppato un motore di automazione in \textbf{Python} che si interfaccia con l'endpoint JSON-RPC di Zabbix. Utilizzando la libreria \texttt{pandas} per il parsing dei fogli Excel, lo script esegue ciclicamente le seguenti operazioni:
\begin{enumerate}
    \item Autenticazione sul server Zabbix e generazione dell'authtoken.
    \item Controllo di esistenza dell'host tramite l'interrogazione \texttt{host.get}.
    \item Se l'host non esiste, invocazione del metodo \texttt{host.create} per inserire l'apparato mappando l'\texttt{ID Host} come nome visibile univoco, associandolo agli Host Groups di competenza territoriale e applicando i template tecnologici corretti.
    \item Valorizzazione dell'\textit{Host Inventory} con i metadati aziendali (es. \texttt{ID Palo}).
\end{enumerate}

\section{Estrazione dati GIS da PDF e georeferenziazione}
Alcune informazioni relative al posizionamento geografico degli apparati risiedevano esclusivamente in relazioni tecniche in formato PDF. Tramite l'utilizzo della libreria Python \texttt{pdfplumber} e l'applicazione di espressioni regolari specifiche (\textit{regex}), \`e stato implementato un sotto-modulo capace di raschiare in automatico i file PDF, estrarre le coordinate GPS (Latitudine e Longitudine) e caricarle nei campi d'inventario di Zabbix, consentendo cos\`i il popolamento in tempo reale delle mappe geografiche (\textit{Geomap}) presenti sulla dashboard di monitoraggio.

\clearpage

% ====================================================================
% CAPITOLO 3
% ====================================================================
\chapter{Monitoraggio Attivo e Gestione degli Eventi TVCC}
\section{Strutturazione gerarchica degli Host Groups}
Per garantire una gestione multi-tenant e una separazione logica dei permessi e delle viste all'interno del frontend di Zabbix, \`e stata implementata una gerarchia di Host Groups organizzata ad albero, sfruttando i separatori nativi (es. \texttt{SmartCity/Comune\_A/TVCC}, \texttt{SmartCity/Comune\_A/Network}).

\section{Middleware per l'integrazione di eventi Hikvision ISAPI}
Il semplice controllo di presenza in rete (Ping) non copre scenari di guasto in cui la telecamera \`e raggiungibile ma il flusso video \`e compromesso (es. obiettivo oscurato o manomesso). Sfruttando le API \textbf{ISAPI} delle telecamere Hikvision, \`e stato sviluppato uno script in Python operante come servizio continuo, il quale si sottoscrive allo stream di eventi HTTP generato dall'apparato all'indirizzo:
\begin{center}
    \texttt{/ISAPI/Event/notification/alertStream}
\end{center}
Il middleware intercetta i messaggi in formato XML/JSON relativi a eventi quali \textit{Video Tampering} (manomissione video) e \textit{Motion Detection} (rilevamento movimento).

\section{Meccanismo di inoltro asincrono con Zabbix Sender}
Non appena il middleware rileva una notifica ISAPI critica, provvede a inoltrare immediatamente l'evento al server Zabbix in modo asincrono, senza attendere il ciclo di polling standard del server. Questa comunicazione avviene tramite l'utility \textbf{Zabbix Sender}, inviando una stringa formattata contenente l'ID dell'host, la chiave dell'item di tipo \textit{Zabbix Trapper} e il valore associato allo stato di allarme.

\section{Configurazione di Trigger e Discovery Actions}
A fronte della ricezione del dato via trapper, specifici \textit{Trigger} valutano le espressioni logiche per determinare lo stato di severit\`a dell'evento. Sono state inoltre configurate delle \textit{Discovery Actions} automatizzate che, in base al tipo di evento ricevuto e all'Host Group di appartenenza, scatenano l'invio immediato di notifiche mail ai tecnici reperibili o avviano script di remediation automatica.

\clearpage

% ====================================================================
% CAPITOLO 4
% ====================================================================
\chapter{Integrazione Crittografica con il Sistema di Ticketing}
\section{Interfacciamento orientato ai servizi: la piattaforma Interacta}
Al fine di standardizzare i flussi di manutenzione sul campo, gli allarmi generati da Zabbix devono essere storicizzati all'interno del sistema di ticketing aziendale basato sulla piattaforma cloud \textbf{Interacta}. L'integrazione richiede che Zabbix sia in grado di aprire un ticket al manifestarsi di un disservizio e di risolverlo automaticamente quando l'apparato torna online.

\section{Autenticazione sicura tramite JSON Web Token (JWT)}
L'interfacciamento con le API REST di Interacta esige un livello elevato di sicurezza basato su autenticazione \textit{Server-to-Server}. Il protocollo prevede lo scambio di un token \textbf{JWT} firmato digitalmente. Il processo si articola nelle seguenti fasi crittografiche:
\begin{enumerate}
    \item Generazione di una coppia di chiavi asimmetriche \textbf{RSA} a 2048 bit.
    \item Configurazione della chiave pubblica sul portale Interacta.
    \item Sviluppo di una funzione Python integrata all'interno degli alert script di Zabbix che compila il payload del JWT (comprensivo dei claim obbligatori come \texttt{iss}, \texttt{sub}, \texttt{iat}, \texttt{exp}).
    \item Firma del token utilizzando la chiave privata custodita a bordo del server Zabbix tramite l'algoritmo asimmetrico \textbf{RS512} (SHA-512 con RSA).
\end{enumerate}

\section{Sviluppo di un Custom AlertScript (Bash/Python)}
L'orchestrazione della chiamata verso Interacta \`e stata implementata tramite un \textbf{Custom AlertScript} depositato nella directory di sistema di Zabbix. Lo script riceve in ingresso i parametri dell'evento nativi di Zabbix (tramite Macro come \texttt{\{EVENT.ID\}}, \texttt{\{HOST.NAME\}}, \texttt{\{TRIGGER.STATUS\}}) ed esegue la richiesta HTTP POST inviando il payload JSON necessario per l'apertura o la chiusura del ticket di intervento.

\clearpage

% ====================================================================
% CAPITOLO 5
% ====================================================================
\chapter{Analisi dei Dati, KPI e Manutenzione del Sistema}
\section{Esportazione continua dei log e superamento dei limiti di retention}
Per permettere analisi storiche a lungo termine senza sovraccaricare il database relazionale di produzione di Zabbix, \`e stato implementato un meccanismo di esportazione continua degli eventi di trigger in file testuali strutturati in formato \textbf{CSV}. 

\section{Algoritmi Python per il calcolo dei KPI operativi}
\`E stato sviluppato un motore di analisi autonomo in Python che elabora periodicamente i file CSV per estrarre metriche prestazionali fondamentali per l'efficienza aziendale:
\begin{itemize}
    \item \textbf{Mean Time To Resolution (MTTR):} Calcolato misurando la differenza temporale tra il timestamp di attivazione del trigger (\texttt{\{EVENT.DATE\} \{EVENT.TIME\}}) e il timestamp di ripristino (\texttt{\{EVENT.RECOVERY.DATE\} \{EVENT.RECOVERY.TIME\}}).
    \item \textbf{Analisi del Flapping:} Algoritmo statistico volto a individuare host instabili che generano frequenti transizioni di stato in finestre temporali ristrette, indice di problemi di alimentazione o connettivit\`a degradata.
    \item \textbf{Distribuzione dei Fault:} Aggregazione dei guasti per area geografica o per tipologia di dispositivo per evidenziare vulnerabilit\`a infrastrutturali ricorrenti.
\end{itemize}

\section{Automazione dei flussi: Crontab e Mail Server Interno}
L'esecuzione del motore di analisi KPI \`e schedulata a livello di sistema operativo Linux tramite il demone \textbf{Cron}. I report generati vengono impacchettati e spediti automaticamente ai responsabili di rete su base giornaliera/settimanale utilizzando uno script Python che sfrutta il protocollo \textbf{SMTP} connesso al mail server interno aziendale.

\section{Politiche di conservazione e manutenzione dei file di log (Logrotate)}
Dato l'elevato volume di log generato sia dagli script di interfacciamento con Interacta che dalle procedure di esportazione KPI, \`e stata configurata una direttiva dedicata all'interno dell'utility di sistema \textbf{Logrotate} (\texttt{/etc/logrotate.d/zabbix-interacta-script}) per prevenire la saturazione del disco, applicando regole di rotazione settimanale, compressione dei file storici e mantenimento di un massimo di 12 istanze.

\clearpage

% ====================================================================
% CONCLUSIONI
% ====================================================================
\chapter*{Conclusioni e Sviluppi Futuri}
\addcontentsline{toc}{chapter}{Conclusioni e Sviluppi Futuri}
Il sistema di monitoraggio centralizzato sviluppato in questa tesi ha risposto con successo alle esigenze operative di Secoval S.r.l., riducendo drasticamente i tempi di configurazione manuale dei nuovi dispositivi e azzerando i tempi di latenza tra il rilevamento di un disservizio sul territorio e la notifica/apertura del ticket nel sistema aziendale Interacta.

L'approccio basato sull'automazione tramite Zabbix API e sullo sviluppo di codice Python integrato ha dimostrato come la figura dell'ingegnere informatico/digitale possa ottimizzare i processi aziendali combinando competenze di networking, sviluppo software e sicurezza dei dati.

Come sviluppi futuri, si prevede l'estensione del sistema attraverso l'implementazione di algoritmi di manutenzione predittiva basati sull'analisi storica delle serie temporali dei dati raccolti, nonch\`e l'integrazione di sensori IoT di tipo ambientale per ampliare il raggio d'azione del monitoraggio Smart City.

\end{document}